\chapter{Non-directed, two-mode and several networks}\label{ch:networktypes}
\begingroup\raggedright\sloppy

The default model of the previous chapters is a single \emph{directed, one-mode, binary} network: an actor controls its own outgoing ties. This chapter covers the three other network situations that RSiena supports and that \nwsaom{} implements: non-directed relations, two-mode (affiliation) networks, and several dependent networks that evolve together. Valued (weighted) networks and signed networks are not supported: a valued wave is rejected with an explicit error and never silently coerced.

\section{Non-directed (symmetric) networks}\label{sec:symmetric}

In a non-directed relation $x_{ij} = x_{ji}$ always holds. A tie change therefore involves two actors, and RSiena has several \emph{model types} for who decides.

\subsection{When the symmetric model is used}

Like RSiena, \nwsaom{} models a relation as non-directed automatically when all waves are declared undirected (\texttt{nwset, \dots{} undirected}, or \texttt{nwsym}) or when they are directed networks that are all tie-symmetric. The option \opt{symmetric} requests it explicitly and is an error for directed waves that are not tie-symmetric: \nwsaom{} never symmetrizes data for you. Waves that mix undirected networks with non-symmetric directed ones are refused. A tie-symmetric directed network stays directed (with a note) when \opt{reciprocity} is requested or when the non-directed model does not apply (network \opt{outdegreeendow}/\opt{reciprocityendow} effects).

For the examples we symmetrise the Glasgow friendship waves (a tie exists if either pupil named the other) and declare them undirected:

\begin{lstlisting}
nwwebuse glasgow, nwclear
forvalues w = 1/3 {
    nwtomata glasgow`w', mat(G`w')
    mata: st_matrix("__S`w'", ((G`w' + G`w'') :> 0))
    nwset, mat(__S`w') undirected name(sg`w') nooutput
}
nwsaom, waves(sg1 sg2 sg3) outdegree gwesp(0.69) nodematch(smoke1) seed(1) k3(300)
\end{lstlisting}
\begin{lstlisting}
SAOM (Method of Moments, conditional), waves: sg1 sg2 sg3
Actors: 50                             Periods: 2
       outdegree |    -2.3699   .1798485   -13.18   0.000
nodematch_smoke1 |   .2077777   .1840164     1.13   0.259
      gwesp_0.69 |   1.615427   .1604325    10.07   0.000
        rate     |    2.4432     2.0675
\end{lstlisting}
No \opt{symmetric} was needed. The engine is native (\texttt{e(engine)} is \texttt{native}); symmetric models require 100\% native term coverage.

\subsection{Model types: symtype()}

\opt{symtype()} chooses how a tie changes. It takes the name or RSiena's \texttt{modelType} number.

\begin{longtable}{@{}p{2.6cm}p{1.8cm}p{2.2cm}p{7.4cm}@{}}
\toprule
\opt{symtype()} & \texttt{modelType} & RSiena code & mechanism \\
\midrule
\endhead
\texttt{forcing} (default) & 2 & AFORCE & An actor, at the per-actor rate, chooses one tie to create or end (or none) exactly as in the directed model and imposes the change. \\
\texttt{confirmation} & 3 & AAGREE & As forcing, but a new tie is only made if the alter confirms it, with the logistic probability of the alter's utility of the new tie. Ending a tie is unilateral. \\
\texttt{force} & 4 & BFORCE & Pairwise: a random pair meets; the tie changes if the \emph{initiator} wants it. \\
\texttt{agree} & 5 & BAGREE & Pairwise: a new tie needs both to want it; either can end a tie. The alter probability is $\sigma(-|u|)$ for utility $u$, exactly as RSiena. \\
\texttt{joint} & 6 & BJOINT & Pairwise: the tie changes if the \emph{summed} utilities of the two favour the change. \\
\bottomrule
\end{longtable}

Before 2026-10-01 the default was \texttt{joint} and types 2 and 3 did not exist; old scripts that relied on the default now fit a different model. The five types on the same data (conditional estimation, density and \opt{nodematch(smoke1)}, \texttt{seed(1)}):

\begin{center}
\begin{tabular}{lrrrr}
\toprule
\opt{symtype()} & outdegree & nodematch(smoke1) & rate, period 1 & rate, period 2 \\
\midrule
forcing & $-1.556$ & 0.234 & 2.011 & 1.774 \\
confirmation & $-0.875$ & 0.152 & 4.112 & 3.536 \\
force & $-2.899$ & 0.395 & 0.321 & 0.270 \\
agree & $-1.325$ & 0.266 & 0.364 & 0.315 \\
joint & $-1.458$ & 0.209 & 0.320 & 0.274 \\
\bottomrule
\end{tabular}
\end{center}
The coefficients are not comparable across model types, because a different mechanism generates the same data. The substantive conclusion about \texttt{nodematch} should be checked under more than one type when the choice is not obvious from theory.

\paragraph{Scale of the rate.} For \texttt{forcing} and \texttt{confirmation} the rate is the per-actor rate, as in the directed model. For the three pairwise types, RSiena gives each actor the basic rate $\lambda$ and each pair the rate $\lambda^2$, while \nwsaom{} simulates a pairwise ministep as an actor with opportunity rate $\rho$ picking an alter among the other $n-1$ actors. The processes are identical with $\lambda^2 = \rho/(n-1)$, and \texttt{e(rate)} is reported on RSiena's scale (\texttt{e(rate\_actor)} keeps $\rho$). In unconditional estimation \texttt{e(rate)} $=\lambda=\sqrt{\rho/(n-1)}$; in conditional estimation it is the mean time to reach the observed distance at basic rate 1, which is on the scale of $\lambda^2$ (on these data 0.30 and 0.55 for \texttt{joint}; $0.55^2 = 0.30$). \opt{rate0()} is on RSiena's scale too. RSiena's default start for the rate of the pairwise types is on the per-actor scale (5.6 on the Glasgow data) and, from it, unconditional estimation often fails; the references in the help file start at 0.5, so give \opt{rate0(0.5)} if an unconditional pairwise fit stops with ``Unlikely to terminate''.

\subsection{Effects in a non-directed network}

Several directed-network effects are meaningless when every tie is symmetric and are rejected (\opt{reciprocity} under \opt{symmetric} stops with ``reciprocity is not meaningful under symmetric'', return code 198): \opt{reciprocity}, \opt{cycle3}, \opt{inactivity}, \opt{outpopularity}, \opt{ininass}, \opt{inoutass}, \opt{outoutass}, \opt{antiiso}, \opt{isolatepop}, \opt{transrectrip} and \opt{transtrip}. These are constants, exact duplicates of another effect, or effects RSiena does not offer for a non-directed variable. Available: \opt{outdegree}, \opt{indegpopularity}, \opt{outactivity}, \opt{cycle4}, \opt{isolatenet}, \opt{outiso}, \opt{antiiniso}, \opt{antiiniso2}, \opt{inplus3}, \opt{outinass}, \opt{gwesp()}, \opt{transties}, \opt{balance}, and every covariate effect (\opt{nodecov()}, \opt{nodeicov()}, \opt{nodeocov()}, \opt{nodematch()}, \opt{simcov()} and their RSiena-spelling aliases). In addition the following exist for non-directed relations: \opt{degplus}, \opt{degplusc} and \opt{transtriads}; the \opt{ownx()} and \opt{netx()} families, including the primary-setting effects; dyadic covariates and the \opt{dyadx()} effects; the behavior-as-covariate effects of Chapter~\ref{ch:coevolution}; and time-varying covariates (\opt{tvcov()} with \opt{symmetric} ran in a test for this manual, although the \opt{tvcov()} entry of the help file lists \opt{symmetric} among the exclusions).

\subsection{Combinations}

\opt{waves()}, \opt{present()}, \opt{missnet()} and \opt{ratecov()} (network-only) combine with \opt{symmetric}. A symmetric network can co-evolve with a behavior (native plugin with the batch path required):
\begin{lstlisting}
nwsaom, waves(sg1 sg2 sg3) outdegree behavior(alcohol1 alcohol2 alcohol3) ///
    linear avalt symtype(joint) seed(1) k3(300)
\end{lstlisting}
\begin{lstlisting}
   outdegree |   -1.32381   .0697569   -18.98   0.000
  beh_linear |   .3504245   .1282042     2.73   0.006
   beh_avalt |   .2587449   .2190706     1.18   0.238
network_rate |    0.5696     0.5264
\end{lstlisting}
(maximum convergence ratio 0.231).

\subsection{Validation}

For \texttt{forcing}, glasgow waves 1--2 symmetrised, density and density plus \opt{nodematch}, conditional and unconditional, five seeds each: every estimate, including the rate, lies within 0.06 RSiena standard errors of RSiena 1.6.6. For the pairwise types (density, \opt{nodematch}, \opt{ratecov}) all estimates including the rates agree within 0.1 standard errors (mean 0.04), except one conditional covariate-rate coefficient (\texttt{force}, 0.13). The coevolution of a symmetric network with a behavior could be checked against RSiena only partly, because RSiena's own estimation of that model is unstable (three of four seeds stop at \texttt{thetaBound}); validation then relies on a simulation study in which RSiena simulated 40 data sets from known parameters and \nwsaom{} re-estimated them (network rate 0.612 for a true 0.6, behavior rate 1.60 for 1.5, outdegree $-1.44$ for $-1.4$, linear median 0.37 for 0.3, avAlt median 1.29 for 1.5, a few replicates blowing up on which RSiena fails too).

\subsection{Limits}\label{sec:sym-limits}

\begin{itemize}
\item \textbf{Several behaviors} with a non-directed network are refused (return code 198).
\item \textbf{Mata engine.} Symmetric co-evolution needs the native batch path; there is no Mata fallback for it.
\item \textbf{Co-evolution with a symmetric network} is only partly cross-validated against RSiena (see above); the avAlt coefficient and the rate scale of that combination are validated by simulation, not by an RSiena fit.
\item \textbf{The RSiena confirmation quirk.} RSiena's \texttt{confirmation} model also computes the confirmation step for the no-change option, which changes nothing but adds to the scores behind the derivative estimate; \nwsaom{} leaves it out.
\item \textbf{No automatic symmetrisation.} Data are never symmetrised by the estimator.
\item \textbf{Rejected effects} listed above; \opt{degplus}, \opt{degplusc} and \opt{transtriads} are only available for non-directed relations, \opt{outpopc} and \opt{inactc} only for directed ones.
\item \textbf{Structural zeros} (\opt{structural()}) are not yet combinable with \opt{symmetric}.
\item \textbf{Degenerate effects.} \texttt{degabsdiffx} and \texttt{degposdiffx} with the behavior as covariate in a non-directed network are available and equal \texttt{siena07}'s statistics; their estimates have large standard errors in both programs. \opt{higher} with a behavior or a changing covariate equals the density statistic in a non-directed network and is not identified (\texttt{siena07} reports a standard error of NA, \nwsaom{} stops).
\end{itemize}

\section{Two-mode networks}\label{sec:twomode}

A two-mode (bipartite) network relates two kinds of nodes, for example persons and the clubs they join. In RSiena's SAOM only the actors of one mode act: an actor changes one tie $x_{ij}$ to an \emph{event} $j$ of the other mode, or does nothing; events never act.

\subsection{Setup}

The network is declared with \texttt{nwset, mat(}$A$\texttt{) bipartite}. By default the \emph{rows} of the affiliation matrix are the actors (stored by nwcommands as mode ``2''); \opt{actormode(1|2)} swaps the roles. Covariates are variables on the \emph{node} rows of the current dataset (one row per node after \texttt{nwset}); the variable \texttt{\_nwmode} tells which mode a row belongs to, and a covariate is read only on the rows of its own mode, centred within that mode. A fully runnable data set (simulated; $40$ actors by $20$ events, three waves) is:

\begin{lstlisting}
clear
nwclear
mata:
mata set matastrict off
void tm_data() {
    na = 40
    ne = 20
    rseed(2026)
    pop = 0.04 :+ 0.25 * runiform(1, ne):^2
    X = (runiform(na, ne) :< pop)
    for (w=1; w<=3; w++) {
        st_matrix("__T" + strofreal(w), X)
        drop = (runiform(na, ne) :< 0.25) :* X
        add  = (runiform(na, ne) :< 0.5 * pop) :* (1 :- X)
        X = X - drop + add
    }
    st_matrix("__CA", rnormal(na, 1, 0, 1))
    st_matrix("__CB", rnormal(ne, 1, 0, 1))
}
tm_data()
mata: CAm = st_matrix("__CA"); CBm = st_matrix("__CB")
end
forvalues w = 1/3 {
    nwset, mat(__T`w') bipartite name(tm`w') nooutput
}
gen double age = .
gen double size = .
mata: md = st_sdata(., "_nwmode"); ia = selectindex(md :== "2"); ie = selectindex(md :== "1")
mata: st_store(ia, "age", CAm); st_store(ie, "size", CBm)
\end{lstlisting}
(\texttt{age} is an actor covariate, \texttt{size} an event covariate.)

\subsection{A model}

\begin{lstlisting}
nwsaom, waves(tm1 tm2 tm3) outdegree cycle4 nodeocov(age) nodeicov(size) seed(1) k3(300)
\end{lstlisting}
\begin{lstlisting}
(two-mode network: 40 actors (mode 2) choose among 20 events)
SAOM (Method of Moments, conditional), waves: tm1 tm2 tm3
Two-mode network: 40 actors choose among 20 events
    outdegree |  -.8976565   .1427513    -6.29   0.000
nodeicov_size |  -.1302557   .0766979    -1.70   0.089
 nodeocov_age |   .0166527   .0868483     0.19   0.848
       cycle4 |    .054978   .0357529     1.54   0.124
        rate  |    2.0851     2.0837
\end{lstlisting}
\texttt{e(twomode)} is 1, and \texttt{e(n\_actors)} and \texttt{e(n\_events)} give the mode sizes (40 and 20). The header line ``Actors: 60'' counts all nodes. With \opt{actormode(1)} the roles swap (20 actors choose among 40 events). In this implementation the network is embedded in a directed one-mode graph on all nodes in which only actor-to-event dyads are free (all others are structural zeros), so the one-mode sampler, plugin and statistics are used unchanged, and the rate is a per-actor rate. Several waves (\opt{waves()}) are supported.

\subsection{Effects available in a two-mode network}

\begin{longtable}{@{}p{4.2cm}p{3.2cm}p{6.4cm}@{}}
\toprule
option & RSiena effect & what it does \\
\midrule
\endhead
\opt{outdegree} & \texttt{density} & Baseline. \\
\opt{indegpopularity} & \texttt{inPopSqrt} & Events with higher degree attract more actors. \\
\opt{outactivity} & \texttt{outAct} & Actors with higher degree make further ties. \\
\opt{cycle4} & \texttt{cycle4} & Four-cycles (two actors sharing two events). \\
\opt{outiso}, \opt{antiiniso}, \opt{antiiniso2}, \opt{inplus3} & \texttt{outIso}, \texttt{antiInIso}, \texttt{antiInIso2}, \texttt{in3Plus} & Isolate and threshold-degree effects. \\
\opt{gwdspfb()}, \opt{avdeg()}, \opt{outtrunc()}, \opt{outmore()}, \opt{outthreshold()} & \texttt{gwdspFB}, \texttt{avDeg}, \texttt{outTrunc}, \texttt{outMore}, \texttt{outThreshold} & Further structural effects. \\
\opt{ownx()} ego-degree effects and \texttt{outinass} & \texttt{outInAss} etc. & Own-network effects without a covariate. \\
\opt{nodeocov(var)} & \texttt{egoX} & Actor covariate, a value per actor. \\
\opt{nodeicov(var)} & \texttt{altX} & Event covariate, a value per event. \\
\opt{egosqx(var)}, \opt{altsqx(var)} & \texttt{egoSqX}, \texttt{altSqX} & Squared covariate effects, actor side / event side. \\
\opt{netx(}\emph{effect}\opt{:}\emph{var}\opt{)} with \emph{effect} one of \texttt{altindist2 totindist2 simegoindist2 sameegoindist2 crossxoutact samexcycle4 avgroupegox totgroupegox} & \texttt{altInDist2}, \texttt{totInDist2}, \texttt{simEgoInDist2}, \texttt{sameEgoInDist2}, \texttt{crossXOutAct}, \texttt{sameXCycle4}, \texttt{avGroupEgoX}, \texttt{totGroupEgoX} & Covariate effects on the \emph{actor} mode. \\
\opt{netx(sameinxcycle4:}\emph{var}\opt{)}, \opt{netx(homxoutact2:}\emph{var}\opt{)} & \texttt{sameInXCycle4}, \texttt{homXOutAct2} & Covariate effects on the \emph{event} mode. \texttt{homxoutact2} needs \opt{xparam(3)} for RSiena's default. \\
\opt{tvcov()} covariates & \texttt{varCovar} & Work in all of these; \texttt{avgroupegox}/\texttt{totgroupegox} only with \opt{tvcov()}, as in RSiena. \\
\bottomrule
\end{longtable}
Coefficients carry the option spelling plus the variable (for example \texttt{nodeicov\_size}, \texttt{egosqx\_age}).

\subsection{A co-evolving behavior of the actors}

An actor behavior may co-evolve with a two-mode network. The events carry no behavior; internally they are generated as \emph{absent actors} (never acting, but remaining among the alters), and their behavior rows are replaced by the actors' mean, so the mean, the range and the centred shape statistics are those of the actors. The network effects that read the behavior are \opt{behegox} and \opt{netbeh(}\texttt{egosqx altindist2 totindist2 simegoindist2 sameegoindist2 samexcycle4 crossxoutact avgroupegox totgroupegox}\opt{)}; these nine are the only \opt{netbeh()} effects of a two-mode network (other tokens used to run on dummy values and are now refused). The events mirror the behavior of the actors, as in RSiena, so there must be at least as many actors as events for the last three.

\begin{lstlisting}
gen double z1 = 0
gen double z2 = 0
gen double z3 = 0
mata: st_store(ia, ("z1","z2","z3"), BAm)     // behavior values of the 40 actors only
nwsaom, waves(tm1 tm2 tm3) outdegree behavior(z1 z2 z3) linear behegox seed(1) k3(300)
\end{lstlisting}
(\texttt{BAm} is an actor-by-wave matrix of integer values 0--3 built in the same Mata block as above.)
\begin{lstlisting}
   outdegree |  -.7255881   .0910187    -7.97   0.000
     behegox |    .079329   .0891253     0.89   0.373
  beh_linear |  -.0497332   .2748901    -0.18   0.856
network_rate |    2.0487     2.0639
behavior_rate|    0.3717     0.7549
\end{lstlisting}
\textbf{Gotcha:} the behavior variable must have non-missing values on the event rows (set them to 0 as above). The same model with Stata missing values (\texttt{.}) on the event rows stopped with the \texttt{thetaBound} divergence error (return code 498), whereas filling the rows with 0 converged.

\subsection{Further options in a two-mode network}
\begin{itemize}
\item \opt{dyadcov()} takes a matrix with one row per actor and one column per event; it is centred by the mean of all its cells (RSiena's bipartite dyadic covariate). \opt{dyadcovrecip()} does not exist in a two-mode network.
\item \opt{interact()} is allowed (for example \texttt{egox} with \texttt{altx}). The interactions of \texttt{egoX} with the density and of the density with \texttt{altX} are degenerate in RSiena itself.
\item The rate may depend on covariates and on the degree: \opt{ratecov()} with constant and changing covariates (the names of \opt{tvcov()}) and with the behavior of the actors as the covariate (this needs \opt{behavior()} and \opt{waves()}), and the structural rate effects \opt{outrate}, \opt{outrateinv} and \opt{outratelog} (Chapter~\ref{ch:rates}). \texttt{inRateInv} and \texttt{inRateLog} do not exist for two-mode networks, because the receivers of the ties are not the actors.
\item A two-mode network can be one of the two networks of \texttt{nwsaom multinet} (Section~\ref{sec:multinet}).
\item Speed: the four-cycle statistic is accumulated for the whole row of alternatives in one sweep, which made two-mode fits with \opt{cycle4} 2.5 times as fast as before (Chapter~\ref{ch:benchmark}).
\end{itemize}

\subsection{Validation}

On a simulated $40\times20$, three-wave network (simulated from known parameters with RSiena itself, since RSiena ships no bipartite example data) the observed statistics equal RSiena's targets (density 529, cycle4 12990, inPopSqrt 2399.29, outAct 4359, egoX 103.142, altX 71.484) and the estimates agree (density $-2.37$ to $-2.50$ against $-2.42$ to $-2.53$; cycle4 0.094 against 0.09--0.096; inPopSqrt 0.57 against 0.57--0.59; rates 4.49 and 3.03 against 4.52 and 3.05). Covariate effects were validated in the same way on a second simulated data set. A two-mode co-evolution with behavior has the base model and six effects checked against \texttt{siena07} (rates 4.31/3.44 against 4.32/3.42).

\subsection{Limits}\label{sec:twomode-limits}

\begin{itemize}
\item \textbf{Effects that do not exist.} \opt{reciprocity}, triadic effects, \opt{balance}, \opt{gwesp} (\opt{reciprocity} stops with return code 198). These have no meaning in a bipartite graph; the RSiena effect list for bipartite variables does not contain them either. Likewise \opt{nodecov()} is refused in a two-mode network (return code 198); use \opt{nodeocov()}/\opt{nodeicov()}.
\item \textbf{Not available in a two-mode network} (return code 198): \opt{endow()}/\opt{creation()}, \opt{missnet()}, a behavior of the events, and a two-mode member of \texttt{nwsaom multinet} together with a behavior. \opt{present()} with all actors present was accepted in a test run; the help file does not list composition change for two-mode networks, so treat it as untested.
\item \textbf{Effects RSiena cannot evaluate either.} \texttt{inPop\_ego}, \texttt{sameXOutAct} and \texttt{diffXOutAct} abort or throw in \texttt{siena07} in a two-mode network and are not offered; \texttt{avGroupEgoX}/\texttt{totGroupEgoX} are not offered with a constant covariate (RSiena reports them as not meaningful); \texttt{inRateInv} and \texttt{inRateLog} stop \texttt{siena07}; \texttt{altInDist2} estimates diverge in \texttt{siena07} too (the statistic equals RSiena's).
\item \textbf{RSiena quirks copied.} \texttt{crossXOutAct} reads the actor covariate at the event's index within its mode for the alter, which needs at least as many actors as events; \texttt{avGroupEgoX} averages over the first $m$ actors ($m$ = number of events). Both behave as in RSiena and may surprise.
\item \textbf{Behavior effects.} With a co-evolving behavior, only \opt{behegox} and the nine \opt{netbeh()} effects above are allowed. Behaviors of events do not exist.
\item \textbf{Only binary affiliation matrices}, as everywhere in \nwsaom{}.
\item \textbf{Missing values in event rows} of a behavior variable (see the gotcha above).
\item \textbf{Wire size.} The structural list of a two-mode network is sent to the plugin as an $n^2$ list, which limits the size of very large bipartite networks (a known limit); the rate start values use $n$ = all nodes.
\item \textbf{Validation base.} No empirical bipartite example exists in RSiena; all validation uses simulated data.
\end{itemize}

\section{Several networks: nwsaom multinet}\label{sec:multinet}

Two networks on the same actors can evolve together, each with its own rate, and the evolution of one may depend on the state of the other (RSiena: two dependent network variables). Each network is directed, non-directed or, for one of them, two-mode (Section~\ref{sec:mn-twomode}). This is a separate subcommand with its own syntax.

\begin{lstlisting}
nwsaom multinet, net1(namelist) net2(namelist) [eff1(effects) eff2(effects)
    cross1(effects) cross2(effects) sym1 sym2 actormode(#) tvcov(groups)
    behavior(namelist) beheff1(effects) beheff2(effects) beheff12(effects)
    beheff21(effects) behint(interactions) missbeh(namelist)
    theta0(numlist) fix(namelist) k0(#) k3(#) firstg(#) seed(#)]
\end{lstlisting}

\opt{net1()} and \opt{net2()} list the waves of the first and second network (same number, at least two; all networks have the same actor set). Estimation is unconditional Method of Moments; every network has a rate per period, estimated jointly with the effects. At every step the network is drawn in proportion to its total rate, then the actor uniformly, then the alter by the network's own objective function (RSiena's \texttt{EpochSimulation}). The engine runs in the plugin only.

\subsection{Within-network and cross-network effects}

\opt{eff1()} and \opt{eff2()} list the within-network effects (default \texttt{outdegree reciprocity}; with \opt{sym1}/\opt{sym2} the default for that network is \texttt{outdegree} only): \texttt{outdegree reciprocity transtrip cycle3 transmedtrip cycle4 indegpop outact}, and with a variable \texttt{egox:}\emph{var}, \texttt{altx:}\emph{var}, \texttt{samex:}\emph{var}, \texttt{simx:}\emph{var}.

\opt{cross1()} lists the effects of network~1 that depend on network~2; \opt{cross2()} the reverse. A cross-network statistic reads the other network at the \emph{start} of the period, the simulation its current state, as RSiena.

\begin{longtable}{@{}p{3.2cm}p{10.6cm}@{}}
\toprule
family & options (lower-case RSiena names) \\
\midrule
\endhead
Basic cross effects & \texttt{crprod crprodrecip crprodmutual from frommutual fromany to toback torecip clxwx clxwx1 clxwx2 mixedinxw} (RSiena \texttt{crprod}, \texttt{crprodRecip}, \texttt{crprodMutual}, \texttt{from}, \texttt{fromMutual}, \texttt{fromAny}, \texttt{to}, \texttt{toBack}, \texttt{toRecip}, \texttt{cl.XWX}, \texttt{cl.XWX1}, \texttt{cl.XWX2}, \texttt{mixedInXW}). \\
Degree families & \texttt{inpopintn inactintn crprodinactintn outpopintn outactintn divoutegointn divinegointn divoutaltintn divinaltintn avdegintn absoutdiffintn both outoutactintn outoutavintn dist2outinactintn outoutdist2actintn outoutdist2avintn ndist2actintn doubleoutact} \\
Closure and mixed triads & \texttt{closure cyclosure sharedin wwxee wwxfr wxxfe wxxer xwxer xwxfe jinmix joutmix gwespffmix gwespbbmix gwespbfmix gwespfbmix gwesprrmix inpopoutw sharedto avalt2mtie avalt2mtot doubleinpop} \\
\bottomrule
\end{longtable}

A number after a colon is RSiena's effect parameter, for example \texttt{from:2}.

\subsection{Cross-network effects with a covariate}

Cross effects that involve a covariate of the actors are written \emph{effect}\texttt{:}\emph{var}[\texttt{:}\emph{\#}], where \emph{var} is a variable in the current dataset and the optional number is RSiena's parameter, for example \texttt{cross1(covnetnet:alcohol1)}. The 34 effects are:
\begin{quote}\small
\texttt{samewwclosure diffwwclosure jumpwwclosure jumpfrom jumpsharedin samewxclosure diffwxclosure jumpwxclosure homwxclosure samexwclosure diffxwclosure covnetnet homcovnetnet contrastcovnetnet alldifcovnetnet covnetnetin homcovnetnetin contrastcovnetnetin alldifcovnetnetin inpopintnx inactintnx outpopintnx outactintnx samexinpopintn samexinactintn samexoutpopintn samexoutactintn altdist2w totdist2w simdist2w altindist2w totindist2w simegodist2w simegoindist2w}
\end{quote}
(RSiena \texttt{sameWWClosure}, \texttt{covNetNet}, \texttt{inPopIntnX}, \texttt{altDist2W}, \dots).

\paragraph{Constants.} The covariate is centred by its mean (RSiena's \texttt{coCovar} convention), which matters for the degree-weighted and distance-two families; the similarity versions use $(x-\bar x)/\text{range}$ with the similarity mean as the parameter. The cross-degree effects use two network constants, taken from the data over all waves: the average degree of both networks and the average squared in-degree of the other network (as RSiena's \texttt{averageInDegree}). No user-supplied constants are needed. RSiena quirks are copied: \texttt{alldifcovnetnetin} is the same function as \texttt{contrastcovnetnetin}; \texttt{DifferentCovariateOutStar} tests the tie alter$\to$$k$ while \texttt{SameCovariateOutStar} tests $k\to$ alter; the average-sameX versions divide by the number of actors with the ego's covariate value.

\subsection{Example}

For the second network we need a relation on the same 50 pupils. We derive an ``advice'' network from friendship (each friendship tie is kept with probability 0.6, plus a few random additions), fixed across waves so that it changes when friendship does. This is only a demonstration device, not a substantive claim.

\begin{lstlisting}
nwwebuse glasgow, nwclear
mata: rseed(99); keep = (runiform(50, 50) :< 0.6); add = (runiform(50, 50) :< 0.02)
forvalues w = 1/3 {
    nwtomata glasgow`w', mat(G`w')
    mata: B = ((G`w' :* keep) + add) :> 0; B = B - diag(diag(B)); st_matrix("__B`w'", B)
    nwset, mat(__B`w') directed name(advice`w') nooutput
}
nwsaom multinet, net1(glasgow1 glasgow2 glasgow3) net2(advice1 advice2 advice3) ///
    eff1(outdegree reciprocity transtrip egox:alcohol1) eff2(outdegree reciprocity) ///
    cross1(crprod covnetnet:alcohol1) cross2(from) seed(1) k3(300)
\end{lstlisting}
\begin{lstlisting}
Several co-evolving networks (Method of Moments), 3 waves, 50 actors
            net1_outdegree |  -2.622782   .1268573   -20.68   0.000
          net1_reciprocity |   2.541702   .2538241    10.01   0.000
            net1_transtrip |   .6400947   .0984247     6.50   0.000
        net1_egox_alcohol1 |   .0086436   .0780513     0.11   0.912
             net1_x_crprod |  -.5927435   .5034914    -1.18   0.239
net1_xc_covnetnet_alcohol1 |   .2379455   .4817869     0.49   0.621
            net2_outdegree |  -1.875265   .1148025   -16.33   0.000
          net2_reciprocity |   1.045423   .3028072     3.45   0.001
               net2_x_from |   .2092609   .3261697     0.64   0.521
Rates (estimated):
           period1  period2
rate_net1    5.632    4.711
rate_net2    1.754    1.749
Overall maximum convergence ratio: 0.242
\end{lstlisting}
Coefficients are named \texttt{net1\_}/\texttt{net2\_} plus the effect; cross effects add \texttt{x\_} (\texttt{net1\_x\_crprod}) or, with a covariate, \texttt{xc\_}\emph{effect}\texttt{\_}\emph{var}. \texttt{e(rates)} has four rows (rate and standard error for each network). On these artificial data \texttt{crprod} is not significant, as expected, and because the advice network is a thinned copy of the friendship network a model with \texttt{crprod} in \emph{both} directions diverged (return code 498) in a trial run: when one network is nearly a subset of the other, the cross effects are close to unidentified.

\subsection{Non-directed members}

\opt{sym1} / \opt{sym2} declares the first / second network non-directed (tie-symmetric or undirected \texttt{nwset} waves; the actor's choice changes the tie in both directions, the forcing model, RSiena's \texttt{AFORCE}, \texttt{symtype} 4 in the plugin). The rate, the distance (two arcs per tie) and the cross statistics are those of the directed engine, and the default within-network effect is \texttt{outdegree} only.
\begin{lstlisting}
nwsaom multinet, net1(sg1 sg2 sg3) net2(advice1 advice2 advice3) sym1 ///
    cross1(crprod from) seed(1) k3(300)
\end{lstlisting}
\begin{lstlisting}
  net1_outdegree |  -1.559588   .1396636   -11.17   0.000
   net1_x_crprod |   .1204944   .4822659     0.25   0.803
     net1_x_from |   .7790062    .278174     2.80   0.005
  net2_outdegree |  -1.841514   .0969473   -19.00   0.000
net2_reciprocity |   1.103243   .2721599     4.05   0.000
Overall maximum convergence ratio: 0.261
\end{lstlisting}
Twenty cross effects were validated against \texttt{siena07} for a non-directed first network with a directed second one (for example \texttt{crprod from fromany to closure clxwx mixedinxw joutmix outpopintn outactintn avdegintn doubleoutact sharedto inpopoutw}).

\paragraph{Earlier subcommand.} The older \texttt{nwsaom multiplex} (two directed networks, exactly two waves each, \texttt{outdegree}, \texttt{reciprocity} and the cross effect \texttt{crprod}/\texttt{crprodb} only) still exists but is deprecated and prints a note; use \texttt{nwsaom multinet}, which generalises it and will replace it.

\subsection{Further cross-network effects}
Besides the families above, the following cross-network effects exist: \texttt{toany}, \texttt{betweenpop} and \texttt{gwespmix} (equal to \texttt{gwespfbmix}); RSiena's \texttt{from.w.ind} as \texttt{fromwindoo} (the other network for both the agreement and the in-degree weights), \texttt{fromwindow} (agreement in the other network, weights from the own) and \texttt{fromwindwo} (agreement in the own, weights from the other). With a \emph{dyadic covariate}, written \emph{effect}\texttt{:}\emph{matrix}[\texttt{:}\emph{\#}] with a Stata matrix, there are \texttt{tou}, \texttt{sharedtou} and \texttt{avaltutie} (RSiena's \texttt{toU}, \texttt{sharedToU}, \texttt{avAltU.2M.tie}), for example \texttt{cross1(sharedtou:W:1)}. As in the current RSiena, \texttt{tou} counts for the ego $i$ and the alter $j$ the actors $q$ with a tie $i\to q$ in the \emph{other} network and a tie $q\to j$ in the own network, each weighted by the covariate value for $(i,q)$. A covariate that changes between waves is declared with \opt{tvcov(a1 a2 a3)}, one variable per period, and used by its stem in the covariate cross effects: \texttt{cross1(samewxclosure:a)}; this does not apply to the three similarity effects.

The cross effects of non-directed pairs are available, too (\texttt{toany gwespmix outoutactintn outoutavintn both avalt2mtie fromwind\ldots{} betweenpop}); a model with a non-directed member was checked against \texttt{siena07} with these effects.

\subsection{A two-mode network as one of the two networks}\label{sec:mn-twomode}
One of the two networks may be two-mode (declared with \texttt{nwset, bipartite}; the network type is detected from the data). Its actors are the nodes of the mode \opt{actormode()} (default 2, the rows of the affiliation matrix); they choose among the events, at a rate of their own in every period, together with the one-mode network, which must have the same actors (the same number; at least 12 nodes in all). The effects of the two-mode network itself are \texttt{outdegree indegpop outact cycle4} (default \texttt{outdegree}), those of the one-mode network the usual ones. The cross-network effects are those of RSiena for such a pair:
\begin{description}[leftmargin=0pt,labelindent=0pt]
\item[Two-mode network depending on the one-mode network] \texttt{inactintn outactintn divoutegointn divinegointn divinaltintn avdegintn outoutactintn outoutavintn dist2outinactintn outoutdist2actintn avalt2mtot to toback torecip sharedto avalt2mtie inpopoutw ndist2actintn toany joutmix gwespfbmix}, and with a covariate of the actors \texttt{samewxclosure:}\emph{var}, \texttt{diffwxclosure}, \texttt{inactintnx}, \texttt{outactintnx}, \texttt{samexinactintn}, \texttt{samexoutactintn} and the distance-two effects \texttt{altdist2w totdist2w simdist2w} (these with at most as many events as actors).
\item[One-mode network depending on the two-mode network] \texttt{outpopintn outactintn divoutegointn divoutaltintn avdegintn outoutactintn outoutavintn dist2outinactintn absoutdiffintn from fromany sharedto avalt2mtie inpopoutw ndist2actintn joutmix gwespmix}; with a covariate of the events (a variable that is missing for the actors) \texttt{altdist2w totdist2w simdist2w} and, with at least as many events as actors, \texttt{outactintnx outpopintnx}.
\item[With a dyadic covariate] \texttt{tou sharedtou avaltutie} (actors by actors for a one-mode other network, actors by events for a two-mode one) and \texttt{fromwindoo}.
\end{description}
\texttt{avalt2mtie} needs the parameter 1 or 2 (give it as \texttt{avalt2mtie:2}). In RSiena some effects index both networks by the same node number, although the alters of a two-mode network are events: \texttt{divinaltintn}, \texttt{joutmix} and \texttt{gwespfbmix} of a two-mode network, and \texttt{outoutactintn} and \texttt{outoutavintn} of a one-mode network that depends on a two-mode one. \nwsaom{} copies this and, as RSiena does, they are available only if the events do not outnumber the actors. The example uses the simulated $40\times20$ data of Section~\ref{sec:twomode} and a one-mode friendship network among the 40 actors that was generated for the example (the first network is the two-mode one).
\begin{lstlisting}
. nwsaom multinet, net1(tm1 tm2 tm3) net2(fr1 fr2 fr3) eff1(outdegree) eff2(outdegree reciprocity) cross1(to) seed(1) k0(100) k3(1000)
Several co-evolving networks (Method of Moments), 3 waves, 40 actors and 20 events (network 1 is two-mode)
-----------------------------------------------------------------------------------
                  | Coefficient  Std. err.      z    P>|z|     [95% conf. interval]
------------------+----------------------------------------------------------------
   net1_outdegree |  -.8501347   .1316513    -6.46   0.000    -1.108167   -.5921028
        net1_x_to |   .2510078   .1819668     1.38   0.168    -.1056406    .6076562
   net2_outdegree |  -1.146603   .0953186   -12.03   0.000    -1.333424   -.9597825
 net2_reciprocity |   .3832283   .2964688     1.29   0.196      -.19784    .9642965
-----------------------------------------------------------------------------------
Rates (estimated):
           period1  period2
rate_net1    2.039    2.068
rate_net2    1.940    1.965
  se_net1    0.285    0.281
  se_net2    0.269    0.245
Overall maximum convergence ratio: 0.098
\end{lstlisting}
Nothing in the cross effect \texttt{to} (actors who are tied to the same friends choose the same events) is expected in these artificial data, and the estimate is not significant.

\subsection{Starting values, fixed effects and tests}
\opt{theta0()} gives the starting values of all effects in the order of the coefficient table (\texttt{net1\_*}, then \texttt{net2\_*}, then the behavior effects), and \opt{fix()} holds the named effects (the coefficient names, or \texttt{\_all}) at those values, as in a single-network fit. After a multinet fit \texttt{estat score} and \texttt{estat timetest} are available (Chapter~\ref{ch:tools}); a fixed cross effect is shown as omitted:
\begin{lstlisting}
. nwsaom multinet, net1(glasgow1 glasgow2 glasgow3) net2(advice1 advice2 advice3) eff1(outdegree reciprocity transtrip) eff2(outdegree reciprocity) cross1(crprod) cross2(from) theta0(-2.6 2.5 0.6 0 -1.9 1 0) fix(net1_x_crprod) seed(1) k3(300)
  net1_outdegree |  -2.644274   .1271252   -20.80   0.000
net1_reciprocity |   2.422096   .2253901    10.75   0.000
  net1_transtrip |   .6082176   .0838987     7.25   0.000
   net1_x_crprod |          0  (omitted)
  net2_outdegree |  -1.878221   .1111325   -16.90   0.000
net2_reciprocity |   1.035088   .3244554     3.19   0.001
     net2_x_from |   .2141975   .2865838     0.75   0.455
\end{lstlisting}
\texttt{estat gof}, \texttt{estat importance} and \texttt{estat simulate} are not available after a multinet fit.

\subsection{Validation}

Observed statistics and estimates of the cross effects agree with \texttt{siena07} on RSiena's s50 friendship data with a second network built from it: the 13 basic cross effects (for example \texttt{crprod} and \texttt{mixedInXW} 1.258 against 1.258; rates 6.1/4.8 and 4.36/4.14 against 6.2/4.8 and 4.36/4.16), and, for the 34 covariate effects, all targets equal and estimates within about 0.05 for the similarity effects. The families in the degree and closure rows were validated against \texttt{siena07} as well.

\subsection{A co-evolving behavior}\label{sec:mn-behavior}

\opt{behavior(\emph{b1 b2 \ldots})} adds one behavior variable (one variable per wave, as in plain \nwsaom{}) that evolves together with the two networks, with its own rate in every period. \opt{beheff1()} lists the behavior effects of the behavior itself and of the first network on it (\texttt{linear quadratic avalt avsim} and the other tokens of the behavior effects of \nwsaom{}; the default is \texttt{linear quadratic avalt}); \opt{beheff2()} lists behavior effects driven by the \emph{second} network, shown as \texttt{beh2\_}\emph{effect}.
\begin{lstlisting}
. nwsaom multinet, net1(glasgow1 glasgow2 glasgow3) net2(nb1 nb2 nb3) behavior(smoke1 smoke2 smoke3) beheff1(linear quadratic avalt) beheff2(avalt) k0(100) k3(1000) seed(1)
  net1_outdegree |  -2.357408   .0979547   -24.07   0.000
net1_reciprocity |   2.820887   .1781864    15.83   0.000
  net2_outdegree |    .220287   .1533902     1.44   0.151
net2_reciprocity |   -.444934   .2799107    -1.59   0.112
      beh_linear |  -1.878621   .3780448    -4.97   0.000
   beh_quadratic |   2.006742   .4216608     4.76   0.000
       beh_avalt |   1.774456    1.18047     1.50   0.133
      beh2_avalt |  -.1019861   1.064681    -0.10   0.924
Rates (estimated):
           period1  period2
rate_net1    5.799    4.521
rate_net2    1.893    2.139
 rate_beh    2.782    2.645
Overall maximum convergence ratio: 0.209
\end{lstlisting}
(\texttt{nb1 nb2 nb3} is a second, randomly generated network in this example.) The behavior can also be a covariate of the network effects, with \texttt{behsim}, \texttt{behegox}, \texttt{behaltx} and \texttt{behegoplusalt} in \opt{eff1()}/\opt{eff2()}, and of the cross-network effects, written \emph{effect}\texttt{:behavior} in \opt{cross1()}/\opt{cross2()}; any of the covariate cross-network effects listed above works this way, for example \texttt{cross1(samewwclosure:behavior covnetnet:behavior)}. In the plugin the behavior statistics are the end behavior on the \emph{starting} network of the period, as in RSiena.

All of these were compared with \texttt{siena07}: the observed statistics are equal (RSiena counts a reciprocated dyad twice and centres the linear and quadratic statistics of a behavior, which changes the targets but not the estimates), and the coefficients and rates agree within simulation noise (at most 0.06 standard errors in the main models), for shape and \texttt{avalt} effects of either network, \texttt{behsim}, \texttt{behegox}, \texttt{behaltx}, and all 34 covariate cross-network effects with the behavior as covariate.

\paragraph{Behavior effects of both networks.} \opt{beheff12()} and \opt{beheff21()} list effects that use the first network as RSiena's first network and the second as its second (and, for \opt{beheff21()}, the networks exchanged). For the actor $i$ with the centred behavior $z_i$ they add $z_i$ times a count over $i$'s alters: \texttt{ffdeg} (alters that $i$ nominates in both networks), \texttt{bbdeg} (in-alters in both), \texttt{fbdeg} (out-alter in the first, in-alter in the second), \texttt{bfdeg} (in-alter in the first, out-alter in the second), \texttt{frdeg} and \texttt{brdeg} (out- and in-alters of the first network that are reciprocated in the second), \texttt{indegmixedpop} and \texttt{outdegmixedpop}. A suffix \texttt{:2} (\texttt{ffdeg:2}) subtracts the degree in the first network (RSiena's parameter 2). Coefficients are named \texttt{beh12\_}\emph{effect} and \texttt{beh21\_}\emph{effect}. All five double-degree effects \texttt{ffdeg}, \texttt{bbdeg}, \texttt{fbdeg}, \texttt{frdeg}, \texttt{brdeg} together are not identified, in RSiena as well.

\paragraph{Covariate, shape and network effects of the behavior.} In \opt{beheff1()} and \opt{beheff2()} a token \emph{family}\texttt{:}\emph{varname} names a covariate effect of the behavior (a numeric variable that is constant over time; at most eight different ones): \texttt{from} (the main effect of the covariate), \texttt{avxalt}, \texttt{avsamexalt}, \texttt{totxalt}, \texttt{totsamexalt}, and the families of \opt{behx()}. The shape effects \texttt{quad\_cc quad\_nc threshold:\#  simallnear:\# simallfar:\# avgroup totgroup} and the further network effects of \opt{behnet2()} are written in the same way. In \opt{beheff1()} the first network drives the effect, in \opt{beheff2()} the second. RSiena's \texttt{constant} effect is confounded with the behavior rate (RSiena returns a standard error of NA), so a fit that includes it stops.

\paragraph{The behavior as a covariate of the network effects.} In \opt{eff1()} and \opt{eff2()}, \texttt{netbeh:}\emph{effect} adds the effect of the centred behavior on tie formation in that network, for example \texttt{eff2(outdegree reciprocity netbeh:absdiffx netbeh:avgroupegox)}. 54 effects of this kind can be used (\texttt{samex diffx higher egosqx altsqx absdiffx egoxaltx samextranstrip \ldots{} avgroupegox totgroupegox}); 53 of them gave the same estimates as \texttt{siena07}, and \texttt{netbeh(indties)} differs from \texttt{siena07} and is refused. The threshold effects take their threshold after a colon.

\paragraph{Endowment, creation, interactions and missing values.} \texttt{linearendow quadraticendow avaltendow avsimendow} in \opt{beheff1()} and \opt{beheff2()} replace the evaluation effect by an endowment and a creation effect (coefficients \texttt{beh\_}\emph{effect}\texttt{\_endow} and \texttt{\_creation}). \opt{behint(}\emph{effect1}\texttt{\#}\emph{effect2}\opt{)} adds the product of two or three behavior effects of the model (RSiena's \texttt{includeInteraction} for behavior effects), each written by its coefficient name, for example \texttt{behint(beh\_quadratic\#beh\_avalt)}; the components must be \texttt{linear}, \texttt{quadratic}, \texttt{avalt} or \texttt{avsim}. \opt{missbeh(m1 m2 \ldots)} names one 0/1 variable per wave that flags missing behavior values; they are replaced by the actor's previous observation, else the next, else the most frequent value of the wave, and do not enter the statistics.

\paragraph{A caution on collinear effects.} A model with \texttt{avalt}, \texttt{avsim} and \texttt{linear} together, or with \texttt{avalt} and \texttt{totalt} of a sparse network, has enormous standard errors in both \nwsaom{} and RSiena on the Glasgow data. The estimates of these effects are correlated at $\pm0.95$ to $\pm1.00$: with a behavior of three categories and 50 actors, most of whom are in the lowest category, the statistics are nearly linear functions of each other and the data can identify only a combination of the coefficients. Use one influence effect at a time.

\subsection{Limits}\label{sec:mn-limits}

\begin{itemize}
\item \textbf{Exactly two networks.} There is no third dependent network, only \opt{net1()} and \opt{net2()}.
\item \textbf{One behavior.} Exactly one behavior can be added. A two-mode member together with a behavior is not available (return code 198).
\item \textbf{Not available inside multinet:} composition change (\opt{present()}, \opt{joinat()}, \opt{leaveat()}), missing ties, \opt{behw()} and \opt{tvbehw()}, \opt{ratecov()}, \opt{endow()}/\opt{creation()} of network effects, \opt{interact()}, \opt{structural()}, \texttt{netbeh(indties)}, and \texttt{estat gof}, \texttt{estat importance} and \texttt{estat simulate}.
\item \textbf{Estimation.} Only \opt{k0()}, \opt{k3()}, \opt{firstg()}, \opt{theta0()}, \opt{fix()} and \opt{seed()} control estimation; conditional estimation does not exist (unconditional only, as in RSiena).
\item \textbf{Refused or not offered.} \texttt{toU} for a one-mode network that depends on a two-mode one is rejected by RSiena; \texttt{avAlt.2M.tie} with RSiena's default parameter 3 throws, so \nwsaom{} offers the parameters 1 and 2; \texttt{outOutDist2AvIntn} has no implementation in RSiena; \texttt{sameWXClosure} for a two-mode network with a behavior is not available.
\item \textbf{Same-actor sets only.} The two networks (or the actors of the two-mode network) must have the same number of nodes.
\end{itemize}
\endgroup
